%
% Template Laporan Skripsi/Thesis 
%
% @author  Andreas Febrian, Lia Sadita 
% @version 1.03
%
% Dokumen ini dibuat berdasarkan standar IEEE dalam membuat class untuk 
% LaTeX dan konfigurasi LaTeX yang digunakan Fahrurrozi Rahman ketika 
% membuat laporan skripsi. Konfigurasi yang lama telah disesuaikan dengan 
% aturan penulisan thesis yang dikeluarkan UI pada tahun 2008.
%

%
% Tipe dokumen adalah report dengan satu kolom. 
%
\documentclass[12pt, a4paper, onecolumn, oneside, final]{report}

% Load konfigurasi LaTeX untuk tipe laporan thesis
\usepackage{_internals/uithesis}
\usepackage{_internals/tikz-uml}
\usepackage{listings}
\usepackage{algorithm}
\usepackage{algpseudocode}
\usepackage{tabularx}
\usepackage{ltablex}
\usepackage[section]{placeins}
\keepXColumns
\usepackage{tocloft}
\usepackage{booktabs}
\usepackage{array}

\setlength{\textfloatsep}{12pt plus 2pt minus 2pt}
\setlength{\floatsep}{10pt plus 2pt minus 2pt}
\setlength{\intextsep}{10pt plus 2pt minus 2pt}

% ===============================
% Daftar Persamaan
% ===============================
\newcommand{\listofequationsname}{DAFTAR PERSAMAAN}
\newlistof{skripsiequations}{loeq}{\listofequationsname}

% Format nomor pada Daftar Persamaan
\renewcommand{\cftskripsiequationspresnum}{\textbf{Persamaan~}}
\renewcommand{\cftskripsiequationsaftersnum}{\normalfont}
\setlength{\cftskripsiequationsnumwidth}{3.7cm}

\newcommand{\eqcaption}[3]{%
\begin{equation}
#1
\label{#2}
\end{equation}
\addcontentsline{loeq}{skripsiequations}{\protect\numberline{\theequation}#3}
}

% Daftar pemenggalan suku kata dan istilah dalam LaTeX
\input{_internals/hype.indonesia}

% Load konfigurasi khusus untuk laporan yang sedang dibuat
%-----------------------------------------------------------------------------%
% Informasi dokumen
%-----------------------------------------------------------------------------%
\var{\judul}{DroneVL: Adaptor VLM dan VLA Agnostik-Model untuk Navigasi UAV Semantik di CoSyS-AirSim}
\Var{\Judul}{\MakeUppercase{\judul}}

\var{\judulInggris}{DroneVL: A Model-Agnostic VLM and VLA Adapter for Semantic UAV Navigation in CoSyS-AirSim}
\Var{\JudulInggris}{\MakeUppercase{\judulInggris}}

\var{\type}{Skripsi}
\Var{\Type}{\MakeUppercase{\type}}

\var{\penulis}{Ryan Adidaru Excel Barnabi}
\Var{\Penulis}{\MakeUppercase{\penulis}}
\var{\npm}{2306266994}

\var{\fakultas}{Teknik}
\Var{\Fakultas}{\MakeUppercase{\fakultas}}
\var{\program}{S1 Teknik Komputer}
\Var{\Program}{\MakeUppercase{\program}}

\Var{\bulan}{[Bulan]}
\Var{\tahun}{2026}
\var{\gelar}{Sarjana Teknik}
\var{\tanggalPengesahan}{[Tanggal pengesahan]}
\var{\tanggalLulus}{[Tanggal lulus]}

\var{\pembimbing}{Alfan Prasekal, S.T., M.Sc. (DIC)}
\var{\kadept}{[Nama Ketua Departemen]}
\var{\kaprodi}{[Nama Ketua Program Studi]}
\var{\pa}{[Nama Pembimbing Akademis]}
\var{\pengujisatu}{[Nama Dosen Penguji 1]}
\var{\pengujidua}{[Nama Dosen Penguji 2]}

\var{\saya}{Penulis}

%-----------------------------------------------------------------------------%
% Judul bab
%-----------------------------------------------------------------------------%
\Var{\kataPengantar}{Kata Pengantar}
\Var{\babSatu}{Pendahuluan}
\Var{\babDua}{Tinjauan Pustaka}
\Var{\babTiga}{Metodologi dan Perancangan Sistem}
\Var{\babEmpat}{Pengujian dan Evaluasi}
\Var{\babLima}{Penutup}

% Daftar istilah yang mungkin perlu ditandai 
\input{istilah}

% Awal bagian penulisan laporan
\begin{document}


% Tambahan CUSTOM 

\definecolor{codegray}{rgb}{0.5,0.5,0.5}
\definecolor{codepurple}{rgb}{0.58,0,0.82}
\definecolor{backcolour}{rgb}{0.95,0.95,0.95}
\definecolor{codeblue}{rgb}{0.0,0.0,0.6}

\lstdefinestyle{mystyle}{
    backgroundcolor=\color{backcolour},
    commentstyle=\color{codegray}\itshape,
    keywordstyle=\color{codeblue}\bfseries,
    stringstyle=\color{codepurple},
    basicstyle=\ttfamily\small,
    numbers=left,
    numberstyle=\tiny\color{codegray},
    stepnumber=1,
    numbersep=8pt,
    breaklines=true,
    breakatwhitespace=true,
    frame=single,
    rulecolor=\color{black},
    tabsize=4,
    captionpos=b,
    showstringspaces=false
}

% Style pseudocode ala artikel ilmiah: caption ringkas di atas, isi kecil,
% spasi rapat, indent tidak terlalu lebar, dan nomor baris kecil.
\floatstyle{ruled}
\restylefloat{algorithm}
\captionsetup[algorithm]{
    font=footnotesize,
    labelfont=bf,
    labelsep=period,
    justification=raggedright,
    singlelinecheck=false,
    skip=4pt
}
\algrenewcommand\algorithmicindent{1.0em}
\algrenewcommand\alglinenumber[1]{\scriptsize\ttfamily\color{codegray}#1}
\algrenewcommand\algorithmicrequire{\textbf{Input:}}
\algrenewcommand\algorithmicensure{\textbf{Output:}}
\newenvironment{compactalgorithmic}[1][1]{%
    \fontsize{1}{1}
    \setstretch{1.0}
    \begin{algorithmic}[#1]
}{%
    \end{algorithmic}
}

% Tambahan CUstom end


%
% Sampul Laporan
\input{_internals/sampul}

%
% Gunakan penomeran romawi
\pagenumbering{roman}

%
% load halaman judul dalam
\addChapter{HALAMAN JUDUL}
\input{_internals/judul_dalam}

%
% setelah bagian ini, halaman dihitung sebagai halaman ke 2
\setcounter{page}{2}

% load halaman orisinalitas 
\chapter*{HALAMAN PERNYATAAN ORISINALITAS}
\thispagestyle{frontmatterNoFooter}
\addcontentsline{toc}{chapter}{HALAMAN PERNYATAAN ORISINALITAS} % masuk daftar isi
%
% Halaman Orisinalitas
%
% @author  Andreas Febrian
% @version 1.01
%

%\chapter*{Halaman Pernyataan Orisinalitas}
\vspace*{2cm}

\begin{center}
\setstretch{2}
	\bo{\type~ini adalah hasil karya saya sendiri, \\
	dan semua sumber baik yang dikutip maupun dirujuk \\
	telah saya nyatakan dengan benar.} \\
	\vspace*{2.6cm}

	\begin{tabular}{l c l}
	\bo{Nama} & : & \bo{\penulis} \\
	\bo{NPM} & : & \bo{\npm} \\
	% Ganti kotak di bawah dengan berkas pindaian tanda tangan, contoh:
	% \raisebox{-0.7\height}{\includegraphics[height=2.4cm]{assets/pics/ttd.png}}
	\bo{Tanda Tangan} & : & \raisebox{-0.7\height}{\framebox[4cm][c]{\rule{0pt}{2.2cm}}} \\
	& & \\
	\bo{Tanggal} & : & \bo{\tanggalPengesahan} \\
	\end{tabular}
\end{center}

\newpage

%
%
% load halaman orisinalitas 
\chapter*{HALAMAN PERNYATAAN PENGGUNAAN KECERDASAN ARTIFISIAL}
\thispagestyle{frontmatterNoFooter}
\addcontentsline{toc}{chapter}{HALAMAN PERNYATAAN PENGGUNAAN KECERDASAN ARTIFISIAL} % masuk daftar isi
%
% Halaman Pernyataan Penggunaan Kecerdasan Artifisial
%
\vspace*{2cm}

\begin{center}
\setstretch{2}
    \bo{Saya menyatakan bahwa dalam penyusunan Skripsi ini, saya menggunakan teknologi Kecerdasan Artifisial melalui platform ChatGPT untuk membantu pemeriksaan dan perbaikan tata bahasa. Perancangan, pelaksanaan penelitian, analisis hasil, penarikan kesimpulan, serta tanggung jawab atas keakuratan, orisinalitas, dan integritas seluruh isi Skripsi tetap berada pada penulis.} \\
    \vspace*{2.6cm}

    \begin{tabular}{l c l}
    \bo{Nama} & : & \bo{\penulis} \\
    \bo{NPM} & : & \bo{\npm} \\
    % Ganti kotak di bawah dengan berkas pindaian tanda tangan, contoh:
    % \raisebox{-0.7\height}{\includegraphics[height=2.4cm]{assets/pics/ttd.png}}
    \bo{Tanda Tangan} & : & \raisebox{-0.7\height}{\framebox[4cm][c]{\rule{0pt}{2.2cm}}} \\
    & & \\
    \bo{Tanggal} & : & \bo{\tanggalPengesahan} \\
    \end{tabular}
\end{center}

\newpage

%
%
\chapter*{HALAMAN PENGESAHAN}
\thispagestyle{frontmatterNoFooter}
\addcontentsline{toc}{chapter}{HALAMAN PENGESAHAN} % masuk daftar isi
\input{src/00-front_matter/03-pengesahan_sidang}
%
%
\chapter*{\kataPengantar}
\thispagestyle{frontmatterNoFooter}
\addcontentsline{toc}{chapter}{\kataPengantar} % masuk daftar isi
%-----------------------------------------------------------------------------%
%\chapter*{Kata Pengantar}
\setstretch{1}
\vspace*{4\baselineskip}
\setstretch{1.5}

Puji dan syukur penulis panjatkan kepada Tuhan Yang Maha Esa atas berkat dan rahmat-Nya sehingga penulis dapat menyelesaikan \type{} dengan judul ``\judul''. Skripsi ini disusun sebagai salah satu syarat untuk memperoleh gelar \gelar{} pada Program Studi \program, Fakultas \fakultas, Universitas Indonesia.

Penulis menyampaikan terima kasih kepada:
\begin{enumerate}
    \item \pembimbing, selaku dosen pembimbing yang telah memberikan arahan, masukan, dan dukungan selama penyusunan skripsi ini;
    \item seluruh dosen dan tenaga kependidikan Program Studi \program{} atas ilmu dan bantuan yang diberikan selama masa studi;
    \item keluarga yang senantiasa memberikan doa dan dukungan; serta
    \item rekan-rekan yang telah memberikan bantuan dan semangat selama proses penelitian dan penulisan skripsi.
\end{enumerate}

Penulis menyadari bahwa skripsi ini masih memiliki keterbatasan. Oleh karena itu, penulis mengharapkan kritik dan saran yang membangun untuk pengembangan penelitian pada masa mendatang. Semoga skripsi ini dapat memberikan manfaat bagi pembaca dan pengembangan penelitian di bidang robotika serta kecerdasan buatan berwujud.

\vspace*{0.1cm}
\begin{flushright}
Depok, \tanggalPengesahan\\[0.1cm]
\vspace*{2cm}
\penulis
\end{flushright}

%
%
\setstretch{1}
\chapter*{Halaman Pernyataan Persetujuan Publikasi Tugas Akhir untuk Kepentingan Akademis}
\thispagestyle{frontmatterNoFooter}
\addcontentsline{toc}{chapter}{HALAMAN PERNYATAAN  PERSETUJUAN PUBLIKASI} % masuk daftar isi
\input{src/00-front_matter/05-persetujuan_publikasi}
%
% 
\fancypagestyle{abstrakstyle}{
    \fancyhf{} % bersihkan header/footer default
    \fancyfoot[R]{\fontsize{10}{12}\selectfont \sffamily \textbf{Universitas Indonesia}}
    \renewcommand{\headrulewidth}{0pt}
    \renewcommand{\footrulewidth}{0pt}
}

\chapter*{ABSTRAK} % judul halaman
\addcontentsline{toc}{chapter}{ABSTRAK} % masuk daftar isi
\thispagestyle{abstrakstyle} % pastikan footer diterapkan
%
% Halaman Abstrak
%
\setstretch{1}
\vspace*{0.2cm}
{
    \singlespacing
    \setlength{\parindent}{0pt}

    \begin{tabular}{@{}l l p{10cm}}
        Nama&: & \penulis \\
        Program Studi&: & \program \\
        Judul&: & \judul \\
        Pembimbing&: & \pembimbing \\
    \end{tabular}

    \bigskip
    \bigskip

\textit{Vision-language model} (VLM) dan \textit{vision-language-action} (VLA) berpotensi digunakan untuk perencanaan semantik serta pengusulan aksi dalam sistem navigasi UAV. Namun, perbedaan prapemrosesan masukan, templat \textit{prompt}, format keluaran, dan ruang aksi menyulitkan penggantian model serta perbandingan yang adil dalam sistem navigasi \textit{closed-loop}. Penelitian ini mengusulkan DroneVL, yaitu adaptor VLM dan VLA agnostik-model untuk navigasi UAV semantik di CoSyS-AirSim yang terintegrasi dengan ROS~2. Eksperimen VLM membandingkan GPT-5.6, Qwen3.6 35B-A3B, dan Gemini Robotics ER 2 dengan varian \textit{Streaming Preview}. Eksperimen portabilitas VLA membandingkan CognitiveDrone, OpenVLA-7B yang diadaptasi untuk UAV, dan OpenPI $\pi_{0.5}$ yang diadaptasi untuk UAV; dua model manipulasi terakhir dilatih menggunakan trajektori CoSyS-AirSim. DroneVL menormalkan citra RGB, data kedalaman, keadaan UAV, dan instruksi misi ke dalam kontrak pengamatan bersama. Setiap respons atau keluaran aksi diurai, dinormalisasi, dan divalidasi menjadi aksi navigasi tingkat tinggi sebelum dijalankan oleh pengendali penerbangan tingkat rendah yang tetap. Eksperimen utama membandingkan model dasar OpenVLA-7B, OpenVLA-7B dengan LoRA yang diadaptasi untuk UAV, dan kondisi LoRA yang dilanjutkan GRPO pada \textit{rollout} simulator. Penelitian ini juga merancang DroneVL Arena, yaitu kerangka \textit{benchmark} dengan episode misi berversi, anotasi semantik, batasan keselamatan, pencatatan lintasan, dan evaluator. Evaluasi mencakup penambatan visual, penalaran spasial, pencarian objek, navigasi semantik, pelaksanaan instruksi multilangkah, keselamatan, efisiensi lintasan, aksi tidak valid, dan latensi. Luaran yang diharapkan adalah lapisan integrasi dan \textit{benchmark} yang dapat direproduksi untuk menganalisis portabilitas dan adaptasi \textit{closed-loop} VLM/VLA pada UAV simulasi.

    \bigskip

    \textbf{Kata kunci:} \textit{vision-language model}, \textit{vision-language-action}, UAV, navigasi semantik, CoSyS-AirSim, ROS~2, LoRA, GRPO, \textit{benchmark}.
}

\newpage

%
%
\chapter*{ABSTRACT} % judul halaman
\addcontentsline{toc}{chapter}{ABSTRACT} % masuk daftar isi
\thispagestyle{abstrakstyle} % pastikan footer diterapkan
%
% Abstract page
%
\setstretch{1}
\vspace*{0.2cm}
{
    \singlespacing
    \setlength{\parindent}{0pt}

    \begin{tabular}{@{}l l p{10cm}}
        Name&: & \penulis \\
        Study Program&: & \program \\
        Title&: & \judulInggris \\
        Supervisor&: & \pembimbing \\
    \end{tabular}

    \bigskip
    \bigskip

Vision-language models (VLMs) and vision-language-action models (VLAs) are promising for semantic planning and action proposal in unmanned aerial vehicle navigation. However, differences in input preparation, prompt templates, response formats, and action spaces make model replacement and fair comparison difficult in a closed-loop navigation system. This research proposes DroneVL, a model-agnostic VLM and VLA adapter for semantic UAV navigation in CoSyS-AirSim integrated with ROS~2. The VLM experiment compares GPT-5.6, Qwen3.6 35B-A3B, and Gemini Robotics ER 2 with the \textit{Streaming Preview} variant. The VLA portability experiment compares CognitiveDrone, UAV-adapted OpenVLA-7B, and UAV-adapted OpenPI $\pi_{0.5}$; the latter two manipulation models are trained on CoSyS-AirSim trajectories. DroneVL normalizes RGB imagery, depth data, UAV state, and mission instructions into a common observation contract. Every response or action output is parsed, normalized, and validated into a high-level navigation action before execution by a fixed low-level flight controller. The primary adaptation experiment compares the base OpenVLA-7B model, UAV-adapted LoRA OpenVLA-7B, and the LoRA condition followed by GRPO on simulator rollouts. This research also designs DroneVL Arena, a benchmark framework with versioned mission episodes, semantic annotations, safety constraints, trajectory logging, and an evaluator. The evaluation covers visual grounding, spatial reasoning, object search, semantic navigation, multi-step instruction following, safety, path efficiency, invalid actions, and latency. The expected output is a reproducible integration layer and benchmark for analysing VLM/VLA portability and closed-loop adaptation in a simulated UAV setting.

    \bigskip

    \textbf{Keywords:} \textit{vision-language model}, \textit{vision-language-action}, UAV, semantic navigation, CoSyS-AirSim, ROS~2, LoRA, GRPO, \textit{benchmark}.
}

\newpage


%
% Daftar isi, gambar, dan tabel
%
\phantomsection
\begingroup
\singlespacing
% --- Aktifkan footer mulai dari sini ---
% Frontmatter style (plain)
\pagestyle{plain}  % Untuk seluruh frontmatter
\tableofcontents
\clearpage
\phantomsection
\listoftables
\clearpage
\phantomsection
\listoffigures
\clearpage
\phantomsection
\addcontentsline{toc}{chapter}{DAFTAR PERSAMAAN}
\listofskripsiequations
\clearpage
\phantomsection
\addcontentsline{toc}{chapter}{\listabbreviationname}
%\listofabbreviation%Jika diperlukan
%\printacronyms
\chapter*{DAFTAR SINGKATAN}
\begin{acronym}[SimpleFlight]
    \acro{AI}{\textit{Artificial Intelligence}}
    \acro{GPS}{\textit{Global Positioning System}}
    \acro{GRPO}{\textit{Group Relative Policy Optimization}}
    \acro{IMU}{\textit{Inertial Measurement Unit}}
    \acro{JSON}{\textit{JavaScript Object Notation}}
    \acro{LiDAR}{\textit{Light Detection and Ranging}}
    \acro{LoRA}{\textit{Low-Rank Adaptation}}
    \acro{PD}{\textit{Proportional-Derivative}}
    \acro{RGB}{\textit{Red, Green, Blue}}
    \acro{ROS}{\textit{Robot Operating System}}
    \acro{UAV}{\textit{Unmanned Aerial Vehicle}}
    \acro{VLM}{\textit{Vision-Language Model}}
\end{acronym}



\clearpage
\endgroup

%
% Gunakan penomeran Arab (1, 2, 3, ...) setelah bagian ini.
%
\pagenumbering{arabic}

%
%
%
% Setelah cover, kata pengantar, abstrak, Daftar Isi
\pagestyle{fancy}
\fancyhf{}
\fancyhead[R]{\thepage} % nomor halaman kanan atas
\fancyfoot[R]{\fontsize{10}{12}\selectfont \sffamily \textbf{Universitas Indonesia}}
\renewcommand{\headrulewidth}{0pt}
\renewcommand{\footrulewidth}{0pt}

\setstretch{1.5} % <-- double spacing
\setlength{\parindent}{1.5em}

%-----------------------------------------------------------------------------%
\chapter{\babSatu}
\label{cha:pendahuluan}
%-----------------------------------------------------------------------------%
Bab ini memaparkan alasan ilmiah dan teknis penelitian DroneVLM. Pembahasan dimulai dari kebutuhan navigasi UAV yang memahami tujuan semantik, kemudian menguraikan masalah integrasi VLM yang heterogen pada sistem navigasi \textit{closed-loop}. Selanjutnya, bab ini merumuskan masalah dan tujuan penelitian, menjelaskan manfaat serta batasan penelitian, dan menempatkan setiap bab sebagai bagian dari alur pembuktian penelitian.

%-----------------------------------------------------------------------------%
\section{Latar Belakang}
\label{sec:latar_belakang}
%-----------------------------------------------------------------------------%
Wahana udara nirawak atau \textit{unmanned aerial vehicle} (UAV) menjalankan misi pada lingkungan yang informasi relevannya dapat berubah selama penerbangan. Pada kondisi tersebut, perencanaan lintasan tidak hanya menentukan rute dari satu koordinat ke koordinat lain, tetapi juga memperbarui keputusan berdasarkan pengamatan lingkungan dan keadaan kendaraan. Tinjauan Vivaldini dkk. menempatkan pengambilan keputusan dari informasi yang diperoleh selama penerbangan sebagai bagian penting dari perencanaan lintasan UAV dalam lingkungan yang tidak pasti \citep{vivaldini2025decision}. Oleh karena itu, navigasi UAV untuk misi semantik memerlukan mekanisme yang dapat menghubungkan pengamatan terkini dengan keputusan gerak berikutnya. % [SRC-B1-01]

Navigasi berbasis tujuan-objek memperluas masalah tersebut karena sasaran tidak selalu tersedia sebagai koordinat yang telah diketahui. UAV perlu mencari objek, mengenali kandidat yang sesuai dengan instruksi, dan memperkirakan lokasi sasaran sebelum membentuk subtujuan geometris yang aman. Ayala dkk. menjelaskan bahwa pencarian dan identifikasi objek merupakan bagian sentral dari navigasi tujuan-objek pada UAV dan berkaitan dengan perencanaan, pelokalan, serta kendali \citep{ayala2024uav}. Dengan demikian, sistem tidak cukup hanya mengoptimalkan lintasan; sistem juga harus memahami apa yang dicari dan bukti visual apa yang mendukung keputusan menuju sasaran tersebut. % [SRC-B1-02]

Instruksi bahasa alami menambah kebutuhan untuk menambatkan entitas bahasa ke lingkungan visual. Agen harus menghubungkan objek, atribut, dan relasi seperti kiri, kanan, depan, atau di antara dengan pengamatan kamera; agen juga harus mengetahui subinstruksi yang telah diselesaikan dan subinstruksi yang perlu dilakukan berikutnya. Penelitian navigasi visi-bahasa menunjukkan bahwa penambatan landmark serta relasi spasial ke modalitas visual penting bagi pelaksanaan instruksi, sedangkan pemantauan kemajuan membantu agen memilih tindakan berikutnya secara konsisten \citep{zhang2022explicit,ma2019selfmonitoring}. Kebutuhan tersebut menjadi lebih menantang pada UAV karena perubahan yaw, ketinggian, dan sudut pandang dapat mengubah bukti visual yang tersedia. % [SRC-B1-03]

Model visi-bahasa atau \textit{vision-language model} (VLM) berpotensi menyediakan penalaran tingkat tinggi yang diperlukan untuk menginterpretasikan pengamatan dan instruksi secara bersama-sama. CLIP menunjukkan bahwa representasi gambar--teks dapat dipelajari dari supervisi bahasa alami dan ditransfer ke konsep visual di luar kelas objek tetap \citep{radford2021clip}. Sementara itu, penelitian pembangkitan keterangan gambar menunjukkan keterhubungan visi komputer dan pemrosesan bahasa alami dalam menghasilkan deskripsi berbasis gambar \citep{vinyals2015show}. Meskipun demikian, keluaran VLM bersifat probabilistik dan dapat berbentuk bahasa bebas, sehingga keluaran tersebut belum dapat diperlakukan sebagai perintah penerbangan tanpa normalisasi dan pemeriksaan lebih lanjut. % [SRC-B1-04]

Masalah integrasi muncul karena setiap \textit{backend} VLM dapat memiliki kebutuhan resolusi gambar, prapemrosesan, templat \textit{prompt}, struktur konteks, dan format respons yang berbeda. Satu model dapat mengembalikan teks bebas, sedangkan model lain dapat mengembalikan pilihan diskret atau objek JSON. Apabila simpul simulator, pengurai respons, dan logika penerbangan ditulis khusus untuk satu \textit{backend}, penggantian model akan mengubah jalur kendali sekaligus. Perbandingan yang dihasilkan menjadi sulit ditafsirkan karena perbedaan kinerja dapat berasal dari perubahan pengurai atau pengendali, bukan dari kemampuan model. Penelitian ini karena itu memerlukan batas antarmuka yang memisahkan inferensi spesifik-model dari simulasi dan kendali yang tetap.

Pemisahan tersebut penting dalam navigasi \textit{closed-loop}, yaitu ketika keputusan saat ini mengubah pengamatan pada langkah berikutnya. Respons VLM yang tidak sesuai skema, ambigu, memakai satuan yang tidak didukung, berasal dari data sensor yang kedaluwarsa, atau mengarahkan UAV ke rintangan tidak boleh langsung dieksekusi. DroneVLM menempatkan VLM sebagai perencana semantik tingkat tinggi, bukan pengendali aktuator. Kontrak pengamatan, pengurai aksi, gerbang keselamatan, pengendali tingkat rendah, mesin keadaan misi, dan pencatat eksperimen dipertahankan sebagai komponen yang terpisah dari \textit{backend}. Rancangan ini membatasi dampak satu respons yang buruk dan memungkinkan sumber kegagalan dicatat secara terpisah.

CoSyS-AirSim digunakan sebagai lingkungan simulasi karena konfigurasi penelitian dapat menyediakan kendaraan multirotor, citra RGB, kedalaman, segmentasi instans, dan lapisan anotasi. Lingkungan tersebut dihubungkan dengan ROS~2 untuk pertukaran citra kamera, odometri, IMU, GPS, data jarak, transformasi objek, perintah kecepatan, dan tujuan posisi lokal. ROS~2 menyediakan abstraksi simpul, topik, layanan, aksi, parameter, serta transformasi yang mendukung integrasi komponen robotika \citep{macenski2022ros2}. Dalam penelitian ini, data anotasi simulator digunakan untuk membangun kebenaran dasar dan mengevaluasi misi, bukan sebagai masukan utama bagi VLM. % [SRC-B1-05]

Repositori dasar penelitian telah menyediakan CoSyS-AirSim, paket jembatan ROS~2, konfigurasi multirotor \texttt{SimpleFlight}, dan pengendali posisi PD sederhana. Namun, repositori tersebut belum menyediakan adaptor VLM, pengelola episode, evaluator arena, maupun \textit{backend} VLM yang siap dibandingkan. Penelitian ini mengisi batas tersebut melalui DroneVLM dan DroneVLM-Arena. DroneVLM menormalkan pengamatan dan respons model, sedangkan DroneVLM-Arena mendefinisikan episode, aturan terminal, anotasi, log, serta evaluator. Dengan pembagian ini, perubahan model tidak memerlukan perubahan pada pengendali atau prosedur evaluasi.

Satu demonstrasi UAV yang mencapai sasaran belum cukup untuk mendukung klaim tentang kemampuan navigasi semantik. Evaluasi harus menggunakan episode dengan pose awal, instruksi, sasaran, batas waktu, anggaran keputusan, dan batasan keselamatan yang ditetapkan terlebih dahulu. Selain tingkat keberhasilan misi, evaluasi perlu mencatat tabrakan, aksi tidak valid, intervensi keselamatan, panjang lintasan, latensi inferensi, dan latensi keputusan ujung-ke-ujung. Benchmark VLM untuk UAV juga menekankan perlunya evaluasi kecerdasan spasial yang mempertimbangkan perspektif serta karakteristik navigasi UAV, bukan hanya jawaban statis dari perspektif manusia \citep{zhang2025skyready}. % [SRC-B1-06]

Berdasarkan kebutuhan tersebut, penelitian ini mengusulkan \textbf{DroneVLM: Adaptor VLM Agnostik-Model untuk Navigasi UAV Semantik di CoSyS-AirSim}. Kontribusi penelitian terdiri atas kontrak pengamatan dan aksi bersama, pemisahan adaptor dari jalur kendali, gerbang keselamatan operasional untuk simulasi, serta arena evaluasi yang dapat direproduksi. Penelitian ini tidak mengusulkan VLM baru dan tidak membuat klaim keselamatan penerbangan di dunia nyata. Hasil yang dilaporkan nantinya dibatasi pada model, adegan, konfigurasi sensor, kondisi lingkungan, dan prosedur evaluasi yang benar-benar diuji dalam DroneVLM-Arena.

%-----------------------------------------------------------------------------%
\section{Rumusan Masalah}
\label{sec:rumusan_masalah}
%-----------------------------------------------------------------------------%
Berdasarkan latar belakang tersebut, rumusan masalah penelitian adalah sebagai berikut.
\begin{enumerate}
    \item Bagaimana merancang dan mengintegrasikan VLM yang heterogen ke dalam satu kerangka navigasi UAV semantik tanpa modifikasi spesifik-model pada tumpukan kendali hilir?
    \item Bagaimana mengevaluasi dan membandingkan VLM pada tugas penambatan visual, penalaran spasial, pencarian objek, navigasi semantik, dan pelaksanaan misi multilangkah dalam lingkungan simulasi yang terstandar?
    \item Sejauh mana adaptasi \textit{supervised fine-tuning} hemat-parameter menggunakan LoRA dan \textit{post-training} berbasis penguatan menggunakan GRPO dapat meningkatkan kinerja navigasi \textit{closed-loop} VLM tujuan umum?
    \item Bagaimana kompromi antara ukuran model, \textit{backend} inferensi, latensi, penggunaan sumber daya komputasi, efisiensi lintasan, keselamatan, dan keberhasilan misi?
    \item Seberapa tangguh perencana UAV berbasis VLM terhadap variasi lingkungan, perumusan \textit{prompt}, modalitas pengamatan, kondisi visual, dan kompleksitas misi; serta bagaimana generalisasinya pada skenario yang belum dilihat?
\end{enumerate}

%-----------------------------------------------------------------------------%
\section{Tujuan Penelitian}
\label{sec:tujuan_penelitian}
%-----------------------------------------------------------------------------%
Selaras dengan rumusan masalah, penelitian ini bertujuan untuk:
\begin{enumerate}
    \item merancang dan mengimplementasikan DroneVLM sebagai adaptor VLM agnostik-model pada CoSyS-AirSim dan ROS~2;
    \item mendefinisikan kontrak pengamatan dan aksi yang memisahkan persiapan serta inferensi spesifik-model dari simulasi, kendali, dan evaluasi;
    \item mengimplementasikan pengurai, penormal aksi, dan gerbang keselamatan untuk memvalidasi aksi tingkat tinggi sebelum dieksekusi;
    \item merancang DroneVLM-Arena yang memuat episode misi berversi, kebenaran dasar semantik, aturan terminal, pencatatan, dan evaluator;
    \item mengevaluasi dan membandingkan VLM terpilih pada penambatan visual, penalaran spasial, pencarian objek, navigasi semantik, dan instruksi multilangkah;
    \item menganalisis keberhasilan misi, keselamatan, efisiensi lintasan, aksi tidak valid, penggunaan sumber daya, dan latensi; serta
    \item mengevaluasi pengaruh adaptasi LoRA dan GRPO terhadap kinerja navigasi \textit{closed-loop} apabila data pelatihan dan sumber daya eksperimen tersedia.
\end{enumerate}

%-----------------------------------------------------------------------------%
\section{Manfaat Penelitian}
\label{sec:manfaat_penelitian}
%-----------------------------------------------------------------------------%
Penelitian ini diharapkan memberikan manfaat sebagai berikut.
\begin{enumerate}
    \item \textbf{Bagi pengembang UAV dan robotika.} DroneVLM menyediakan batas yang jelas antara \textit{backend} VLM, simulator, validasi keselamatan, dan pengendali penerbangan. Pengembang dapat menguji \textit{backend} baru melalui adaptor tanpa menulis ulang jalur ROS~2, pengendali tingkat rendah, atau evaluator.
    \item \textbf{Bagi peneliti \textit{embodied AI}.} DroneVLM-Arena menyediakan rancangan evaluasi \textit{closed-loop} dengan episode dan prosedur pencatatan yang tetap. Rancangan tersebut membantu membedakan kegagalan inferensi, pengurai, keselamatan, pengendali, dan infrastruktur dari sekadar hasil akhir misi.
    \item \textbf{Bagi pengguna akademik.} Penelitian ini menyediakan contoh integrasi VLM--ROS~2 pada UAV simulasi dan dasar untuk penelitian lanjutan tentang perencanaan tangguh, keselamatan operasional, serta \textit{sim-to-real transfer}. Hasil evaluasi dapat digunakan untuk menentukan batas kemampuan sistem pada kondisi yang diuji.
    \item \textbf{Bagi perancang sistem.} Analisis metrik kualitas, keselamatan, efisiensi, sumber daya, dan latensi dapat menjadi dasar pemilihan \textit{backend} sesuai kebutuhan misi. Pemilihan tersebut tidak hanya mempertimbangkan keberhasilan misi, tetapi juga biaya dan risiko operasional pada jalur keputusan.
\end{enumerate}

%-----------------------------------------------------------------------------%
\section{Ruang Lingkup Penelitian}
\label{sec:ruang_lingkup}
%-----------------------------------------------------------------------------%
Ruang lingkup penelitian dibatasi sebagai berikut.
\begin{enumerate}
    \item Evaluasi dilakukan pada CoSyS-AirSim Blocks. Penelitian tidak mencakup penerbangan dunia nyata, pengujian \textit{hardware-in-the-loop}, maupun klaim keselamatan penerbangan nyata.
    \item Kendaraan yang digunakan adalah satu UAV multirotor dalam mode \texttt{Multirotor} dengan \texttt{SimpleFlight}. Pengendali tingkat rendah dibuat tetap pada seluruh perbandingan \textit{backend}.
    \item VLM berperan sebagai perencana semantik tingkat tinggi dan tidak dapat menerbitkan perintah motor atau gaya dorong secara langsung.
    \item Pengamatan utama terdiri atas citra RGB depan, kedalaman terdaftar, keadaan atau odometri UAV, ringkasan kemajuan misi, dan instruksi bahasa alami. LiDAR, fusi banyak kamera, SLAM penuh, dan memori jangka panjang berada di luar konfigurasi utama kecuali sebagai ablasi.
    \item \textit{Backend} utama yang dikunci adalah LLaVA-v1.5-7B, Qwen2-VL-2B-Instruct, dan InternVL2-2B. Revisi model, checksum SHA-256 bobot, kuantisasi, \textit{runtime}, prosesor, \textit{prompt}, serta parameter dekode dicatat sebelum pengujian.
    \item Setiap \textit{backend} menggunakan skema aksi, pengendali, batas gerak, anggaran waktu dan keputusan, geofence, gerbang keselamatan, kebijakan coba-ulang, serta evaluator yang sama. Perbedaan kemampuan bawaan dilaporkan sebagai hasil, tidak dinormalisasi atau disembunyikan.
    \item Benchmark dibatasi pada penambatan visual, penalaran spasial, pencarian objek, navigasi semantik, dan instruksi multilangkah. Manipulasi, koordinasi banyak UAV, pendaratan presisi, dan interaksi fisik tidak dibahas.
    \item Label simulator dan transformasi objek hanya digunakan untuk anotasi, baseline oracle, dan evaluasi; keduanya tidak diberikan kepada VLM pada kondisi utama.
    \item Kesimpulan dibatasi pada adegan, konfigurasi sensor, distribusi bahasa, model, serta kondisi lingkungan yang ditetapkan dan diuji dalam DroneVLM-Arena.
\end{enumerate}

%-----------------------------------------------------------------------------%
\section{Sistematika Penulisan}
\label{sec:sistematika}
%-----------------------------------------------------------------------------%
Laporan ini disusun dalam lima bab yang saling berkaitan.
\begin{enumerate}
    \item \textbf{Bab I Pendahuluan} memaparkan konteks masalah, rumusan masalah, tujuan, manfaat, ruang lingkup, dan alur penulisan penelitian.
    \item \textbf{Bab II Tinjauan Pustaka} menjelaskan teori yang mendasari DroneVLM, penelitian terkait, celah penelitian, dan posisi kontribusi yang diusulkan.
    \item \textbf{Bab III Metodologi dan Perancangan Sistem} menjabarkan desain penelitian, rancangan arsitektur, kontrak antarkomponen, episode evaluasi, metrik, dan protokol pencatatan.
    \item \textbf{Bab IV Pengujian dan Evaluasi} menyajikan konfigurasi eksperimen, hasil, validasi data, dan analisis keberhasilan misi, keselamatan, efisiensi, serta latensi berdasarkan log evaluator.
    \item \textbf{Bab V Penutup} merangkum jawaban atas rumusan masalah berdasarkan bukti Bab IV dan menyajikan saran pengembangan penelitian berikutnya.
\end{enumerate}

%-----------------------------------------------------------------------------%
\chapter{\babDua}
\label{cha:studiliteratur}
%-----------------------------------------------------------------------------%
Bab ini menyajikan teori dan penelitian terkait yang diperlukan untuk merancang serta mengevaluasi DroneVL. Pembahasan dimulai dari navigasi UAV dan kemampuan model VLM dan VLA, kemudian menjelaskan penambatan visual, penalaran spasial, navigasi \textit{closed-loop}, integrasi CoSyS-AirSim--ROS~2, kontrak aksi, dan evaluasi yang dapat direproduksi. Bagian akhir membandingkan pendekatan terdahulu dengan kontribusi DroneVL untuk mengidentifikasi celah penelitian yang diuji.

%-----------------------------------------------------------------------------%
\section{Landasan Teoretis}
\label{sec:landasan_teori}
%-----------------------------------------------------------------------------%

\subsection{UAV dan Navigasi Otonom}
\label{subsec:uav_navigasi}
UAV adalah kendaraan udara yang dapat beroperasi tanpa pilot di dalam kendaraan. Pada multirotor, gaya dorong beberapa rotor menghasilkan gerak translasi dan rotasi, sedangkan tumpukan navigasi mengestimasi keadaan kendaraan, membentuk sasaran, dan meneruskan target posisi, kecepatan, ketinggian, atau yaw kepada pengendali penerbangan. Keputusan selama penerbangan perlu mempertimbangkan informasi yang diperoleh dari lingkungan karena ketidakpastian dapat mengubah lintasan yang layak ditempuh \citep{vivaldini2025decision}. Pada navigasi tujuan-objek, kebutuhan tersebut mencakup pencarian, pengenalan, dan pelokalan objek sasaran sebelum kendaraan dapat membentuk rute menuju sasaran \citep{ayala2024uav}. % [SRC-B2-01]

Navigasi geometris umumnya menyatakan tujuan melalui koordinat, \textit{waypoint}, peta, atau representasi spasial lain yang dapat langsung dipakai perencana. Navigasi semantik menyatakan tujuan sebagai objek, atribut, relasi, atau instruksi; sistem harus terlebih dahulu menambatkan deskripsi tersebut, kemudian membentuk satu atau lebih subtujuan geometris. Kedua pendekatan tidak saling meniadakan: DroneVL menggunakan penalaran semantik untuk menentukan sasaran atau relasi yang relevan, lalu menggunakan kendali dan perencanaan geometris untuk mencapai subtujuan secara aman. Perbedaan ini penting karena objek dapat tertutup, berada di luar sudut pandang awal, atau memiliki beberapa kandidat yang serupa. % [SRC-B2-02]

\subsection{\textit{Vision-Language Model}}
\label{subsec:vlm}
\textit{Vision-language model} (VLM) menghubungkan representasi visual dan bahasa melalui penyandi visual, model bahasa, serta mekanisme penyelarasan lintasmodal. CLIP memperlihatkan bahwa supervisi bahasa alami dapat menghasilkan representasi gambar--teks yang dapat ditransfer ke tugas visual, sehingga pendekatan tersebut tidak terbatas pada daftar kelas objek yang tetap \citep{radford2021clip}. Penelitian \textit{Show and Tell} juga menunjukkan bahwa arsitektur saraf dapat mengaitkan visi komputer dengan pemrosesan bahasa alami untuk menghasilkan deskripsi gambar \citep{vinyals2015show}. Landasan ini mendukung penggunaan VLM sebagai komponen penalaran semantik, tetapi tidak menjadikan keluaran bahasa sebagai bukti fisik atau perintah penerbangan yang aman. % [SRC-B2-03]

Keluaran VLM bersifat probabilistik dan peka terhadap resolusi, pemotongan gambar, urutan konteks, templat \textit{prompt}, dan parameter dekode. Model dapat salah mengenali objek, menghalusinasi objek yang tidak terlihat, atau menghasilkan respons yang tidak sesuai dengan tindakan yang tersedia. Oleh karena itu, DroneVL memperlakukan VLM sebagai komponen inferensi dengan kontrak masukan dan keluaran yang ketat. Kontrak tersebut membatasi perbedaan antarmodel pada adaptor, sementara simulator, pengendali, evaluator, dan kebijakan keselamatan tetap menggunakan representasi bersama. Dengan pemisahan ini, kesalahan respons dapat ditelusuri tanpa menyamakan respons bahasa model dengan keadaan sebenarnya dari lingkungan.

\subsection{\textit{Vision-Language-Action}}
\label{subsec:vla}
Model \textit{vision-language-action} (VLA) menghubungkan pengamatan visual dan instruksi bahasa dengan keluaran aksi kebijakan. OpenVLA-7B adalah VLA terbuka yang menghasilkan token aksi diskret dan menyediakan prosedur \textit{fine-tuning} LoRA untuk tugas serta perwujudan robot baru \citep{kim2024openvla}. OpenPI menyediakan keluarga $\pi_0$, $\pi_0$-FAST, dan $\pi_{0.5}$; implementasi $\pi_{0.5}$ saat ini menggunakan kepala \textit{flow-matching} untuk menghasilkan potongan aksi kontinu \citep{physicalintelligence2025openpi}. CognitiveDrone berbeda karena dirancang khusus untuk UAV, dilatih pada trajektori penerbangan simulasi, dan menghasilkan perintah aksi empat-dimensi dari citra orang-pertama serta instruksi \citep{lykov2025cognitivedrone}.

Keluaran aksi VLA selalu terikat pada perwujudan dan data pelatihan asalnya. Karena OpenVLA-7B dan OpenPI $\pi_{0.5}$ ditujukan untuk manipulasi robot, aksi sendi atau potongan aksi asal tidak dapat dipetakan langsung menjadi kendali UAV. DroneVL hanya mengevaluasi versi yang telah diadaptasi terhadap skema aksi UAV CoSyS-AirSim melalui trajektori berpasangan RGB--instruksi--keadaan UAV--aksi UAV. Setiap skema menentukan satuan, kerangka koordinat, horizon aksi, dan prosedur dekode sebelum keluaran dapat menjadi kandidat aksi kanonis (\textit{canonical action}). Batas ini membedakan adaptasi kebijakan yang dapat diuji dari klaim transfer lintas-perwujudan \textit{zero-shot}.

\subsection{Pemilihan Model untuk Eksperimen}
\label{subsec:pemilihan_model}
Pemilihan model pada DroneVL bersifat bertujuan, bukan peringkat kemampuan umum. Model dipilih apabila dapat menerima pengamatan visual dan instruksi, menyediakan antarmuka keluaran yang dapat dinormalisasi ke kontrak DroneVL, serta mewakili variasi yang relevan bagi pertanyaan penelitian: akses layanan tertutup atau model yang dapat dikonfigurasi, keluaran terstruktur atau pemanggilan fungsi, penalaran berwujud, dan kebijakan aksi yang terikat pada perwujudan. Kapabilitas yang dideklarasikan penyedia dicatat sebagai konfigurasi awal; keunggulan pada penambatan visual, keselamatan, efisiensi lintasan, dan latensi hanya disimpulkan dari eksperimen pada Bab~III dan Bab~IV.

\subsubsection{Kelompok VLM}
GPT-5.6 dipilih sebagai pembanding VLM berbasis API untuk penalaran multimodal tingkat tinggi. Dokumentasi OpenAI menyatakan bahwa alias \texttt{gpt-5.6} merujuk ke GPT-5.6 Sol dan mendukung masukan citra, keluaran terstruktur, serta pemanggilan fungsi \citep{openai2026gpt56}. Karakteristik ini memungkinkan keluaran perencana dikirim dalam format yang dapat diperiksa oleh pengurai DroneVL. Model ini tidak digunakan sebagai pengendali penerbangan, melainkan sebagai sumber keputusan semantik yang seluruh keluarannya tetap melewati normalisasi dan gerbang keselamatan.

Qwen3.6 35B-A3B dipilih sebagai VLM \textit{mixture-of-experts} (MoE) multimodal yang konfigurasi \textit{runtime} dan prapemrosesannya dapat didokumentasikan secara eksplisit pada manifest eksperimen. Varian yang digunakan adalah \texttt{qwen3.6:35b-a3b} melalui Ollama; halaman model mencatat bobot 35,5B, arsitektur \texttt{qwen35moe}, kuantisasi Q4\_K\_M, dan proyektor CLIP \citep{ollama2026qwen36}. Dokumentasi Qwen juga mencantumkan dukungan masukan teks, citra, dan video, pemanggilan fungsi, serta keluaran terstruktur untuk varian 35B-A3B \citep{qwen2026visionmodels}. Karena itu, Qwen3.6 35B-A3B memungkinkan pengujian apakah kontrak pengamatan dan aksi DroneVL tetap berlaku ketika model, cara penyajian citra, dan lingkungan inferensi berbeda dari layanan API tertutup. Tag model, checksum, kuantisasi, dan perangkat keras yang benar-benar digunakan dikunci sebelum pengujian.

Gemini Robotics ER 2 dengan varian \textit{Streaming Preview} dipilih karena dirancang untuk penalaran berwujud, termasuk penalaran spasial, pemahaman kemajuan video, dan orkestrasi tugas multilangkah. Varian ini juga dioptimalkan untuk interaksi dua arah berlatensi rendah melalui Live API dan mendukung pemanggilan fungsi \citep{google2026geminirobotics}. Dengan demikian, Gemini menguji apakah adaptor DroneVL dapat menangani respons \textit{streaming} dari VLM yang berorientasi robotika tanpa mengubah pengendali UAV. Fokus evaluasinya tetap pada kualitas keluaran yang berhasil dinormalisasi dan dieksekusi secara aman dalam simulator, bukan pada klaim kemampuan robotika penyedia.

Ketiga VLM tersebut membentuk pembandingan antarmuka yang saling melengkapi: GPT-5.6 sebagai layanan API multimodal dengan keluaran terstruktur, Qwen3.6 35B-A3B sebagai VLM MoE yang dijalankan secara lokal dengan konfigurasi inferensi yang dapat dikendalikan dan dicatat, serta Gemini sebagai VLM berorientasi robotika dengan jalur \textit{streaming}. Perbedaan akses, modalitas, dan antarmuka sengaja tidak dihapus; seluruhnya dilaporkan bersama hasil agar perbandingan tidak menganggap ketiga model memiliki kondisi operasi yang identik.

\subsubsection{Kelompok VLA}
CognitiveDrone dipilih sebagai titik acuan VLA yang sejak awal ditujukan bagi UAV. Model tersebut menggunakan pengamatan orang-pertama dan instruksi untuk menghasilkan aksi empat-dimensi pada trajektori penerbangan simulasi \citep{lykov2025cognitivedrone}. Perannya adalah menguji adaptor pada kebijakan yang sudah dekat dengan perwujudan UAV, sambil tetap mewajibkan deklarasi skema aksi, satuan, kerangka koordinat, dan horizon keluaran.

OpenVLA-7B dipilih sebagai VLA utama untuk eksperimen adaptasi karena keluaran aksinya berupa token diskret dan implementasinya menyediakan prosedur \textit{fine-tuning} LoRA untuk tugas serta perwujudan robot baru \citep{kim2024openvla}. Sifat tersebut mendukung tiga kondisi penelitian: model dasar dengan tokenisasi aksi UAV (M0), adaptasi LoRA pada trajektori CoSyS-AirSim (M1), dan M1 yang dilanjutkan GRPO (M2). Dengan kata lain, OpenVLA-7B bukan hanya pembanding portabilitas, tetapi objek utama untuk menguji pengaruh LoRA dan GRPO pada kebijakan UAV yang menggunakan kontrak aksi yang sama.

OpenPI $\pi_{0.5}$ dipilih untuk menguji batas portabilitas adaptor terhadap VLA dengan representasi aksi yang berbeda. OpenPI menyediakan kebijakan yang menghasilkan potongan aksi kontinu melalui kepala \textit{flow-matching} \citep{physicalintelligence2025openpi}; setelah adaptasi UAV, keluarannya harus diterjemahkan melalui \textit{NativeActionEnvelope} sebelum menjadi kandidat aksi kanonis. Model ini tidak dimasukkan ke eksperimen GRPO, karena keluaran kontinu tersebut tidak serta-merta menyediakan rasio probabilitas kebijakan yang diperlukan oleh prosedur GRPO. Perannya adalah menguji kontrak dan pengaman DroneVL pada perbedaan representasi aksi, bukan membandingkan metode optimisasi yang tidak sepadan.

Secara bersama, kelompok VLA mencakup kebijakan yang dirancang untuk UAV, VLA bertoken diskret yang menjadi objek adaptasi utama, dan VLA dengan aksi kontinu. Susunan ini memungkinkan penelitian membedakan dua pertanyaan: apakah adaptor dapat memportasikan model dengan asumsi perwujudan serta format aksi yang berbeda, dan apakah LoRA+GRPO pada OpenVLA-7B yang telah diadaptasi meningkatkan kinerja terhadap kondisi dasarnya. Kedua pertanyaan dilaporkan terpisah agar hasil portabilitas tidak ditafsirkan sebagai efek LoRA atau GRPO.

\subsection{\textit{Low-Rank Adaptation} dan \textit{Group Relative Policy Optimization}}
\label{subsec:lora_grpo}
\textit{Low-rank adaptation} (LoRA) adalah pendekatan \textit{parameter-efficient fine-tuning} yang membekukan bobot model pralatih dan menambahkan matriks berperingkat rendah yang dapat dilatih pada lapisan transformator \citep{hu2021lora}. Pendekatan tersebut relevan bagi model terbuka karena adaptasi domain tidak memerlukan pembaruan seluruh parameter. QLoRA menggabungkan LoRA dengan kuantisasi empat-bit pada model pralatih yang dibekukan untuk mengurangi kebutuhan memori selama \textit{fine-tuning} \citep{dettmers2023qlora}. Dalam DroneVL, LoRA dan QLoRA hanya diterapkan pada OpenVLA-7B sebagai VLA utama; pendekatan ini tidak diterapkan pada VLM API tertutup maupun pada Qwen3.6 35B-A3B agar efek adaptasi kebijakan aksi UAV dapat diisolasi. % [SRC-B2-10]

\textit{Group Relative Policy Optimization} (GRPO) diperkenalkan sebagai varian \textit{proximal policy optimization} yang menggunakan perbandingan hasil dalam satu kelompok respons dan ditujukan untuk menurunkan kebutuhan memori dibandingkan pendekatan dengan \textit{critic} terpisah \citep{shao2024deepseekmath}. Dalam ranah multimodal, Vision-R1 menerapkan GRPO bersama penghargaan format dan data penalaran multimodal, sedangkan VLM-R1 mengeksplorasi pembelajaran penguatan bergaya R1 pada tugas visual dengan \textit{ground truth} yang terdefinisi \citep{huang2025visionr1,shen2025vlmr1}. Pada DroneVL, eksperimen utama LoRA+GRPO diterapkan pada OpenVLA-7B yang telah diberi tokenisasi aksi UAV; keluaran aksi diskretnya memungkinkan probabilitas token aksi dari setiap \textit{rollout} dicatat untuk rasio kebijakan GRPO. Tiga kondisi dibandingkan pada set pengujian yang sama dan telah dibekukan: model dasar OpenVLA-7B (M0), OpenVLA-7B dengan LoRA yang diadaptasi untuk UAV (M1), dan M1 yang dilanjutkan GRPO pada \textit{rollout} CoSyS-AirSim (M2). Fungsi penghargaan, data pelatihan, anggaran \textit{rollout}, komputasi, dan komponen ablasinya ditetapkan di Bab~III sebelum pelatihan dimulai. OpenPI $\pi_{0.5}$ tetap menjadi eksperimen portabilitas VLA terpisah; metode \textit{flow-matching} kontinu tidak disebut GRPO tanpa estimator probabilitas kebijakan yang tervalidasi. % [SRC-B2-11]

\subsection{Penambatan Visual, Penalaran Spasial, dan Instruksi Multilangkah}
\label{subsec:grounding_spasial}
Penambatan visual adalah pemetaan entitas yang disebut dalam instruksi ke objek atau wilayah pada pengamatan. Proses ini mencakup identifikasi kelas, atribut, dan instans yang benar di antara objek pengalih. Dalam \textit{vision-language navigation}, landmark serta relasi spasial pada instruksi perlu ditambatkan ke modalitas visual agar agen dapat memilih arah yang sesuai \citep{zhang2022explicit}. DroneVL Arena karena itu mengevaluasi dua keluaran yang berbeda: kecocokan instans target yang dideklarasikan terhadap \textit{ground truth}, serta terpenuhinya predikat spasial oleh pose UAV terhadap target. Setiap episode menyatakan kerangka acuan, toleransi geometris, dan waktu evaluasi bagi relasi seperti kiri, kanan, atau depan. % [SRC-B2-04]

Instruksi multilangkah menambahkan kendala urutan pada masalah penambatan. Pencapaian wilayah akhir tidak selalu berarti misi berhasil apabila subtujuan sebelumnya belum dipenuhi. Agen perlu membedakan instruksi yang telah selesai, instruksi yang menjadi fokus berikutnya, arah gerak yang relevan, dan kemajuan menuju sasaran. Kebutuhan tersebut menjadi dasar pendekatan \textit{self-monitoring} pada \textit{vision-language navigation} \citep{ma2019selfmonitoring}. DroneVL menggunakan riwayat tindakan yang terbatas dan ringkasan kemajuan misi sebagai bagian pengamatan kanonis agar keputusan berikutnya dapat mempertimbangkan konteks tanpa membuat masukan model bertumbuh tanpa batas. % [SRC-B2-05]

\subsection{\textit{Embodied Closed-Loop Agent}}
\label{subsec:closed_loop}
\textit{Embodied agent} menjalankan lingkar persepsi--keputusan--aksi--pengamatan. Berbeda dengan evaluasi gambar statis, tindakan pada navigasi \textit{closed-loop} mengubah keadaan dan pengamatan pada langkah berikutnya. Rotasi dapat membuka pandangan baru ke sasaran, sedangkan gerakan maju yang salah dapat mengurangi jarak bebas, menyebabkan tabrakan, atau menghabiskan anggaran misi. Karena itu, evaluasi harus mengukur konsekuensi rangkaian keputusan, bukan hanya ketepatan satu jawaban terhadap satu gambar.

DroneVL menerapkan horizon aksi pendek dan keputusan berulang. Setelah satu aksi atau satu durasi terbatas, sistem memperoleh pengamatan baru, membangun kembali konteks misi, dan meminta keputusan berikutnya dari \textit{backend}. Batas kecepatan dan durasi membatasi paparan serta perpindahan akibat satu keputusan, tetapi tidak menjamin bahwa akibatnya dapat dipulihkan; karena itu gerbang keselamatan tetap memeriksa jarak bebas, geofence, usia pengamatan, dan kondisi penghentian. Pengelola misi menentukan keadaan misi dan kondisi terminal, sedangkan pengendali tingkat rendah tetap sama pada seluruh \textit{backend}. Kesamaan ini mengendalikan perubahan pada jalur hilir, sementara perbedaan modalitas, resolusi, latensi, dan kapabilitas \textit{backend} dicatat sebagai faktor eksperimen, bukan diasumsikan setara.

\subsection{CoSyS-AirSim dan ROS~2}
\label{subsec:cosys_ros}
AirSim menyediakan simulasi fisik dan visual untuk kendaraan otonom \citep{shah2017airsim}, sedangkan CoSyS-AirSim memperluas kerangka tersebut untuk kebutuhan simulasi industri yang lebih kompleks \citep{jansen2023cosysairsim}. CoSyS-AirSim digunakan penelitian ini untuk membangun episode UAV multirotor dengan citra, kedalaman, segmentasi instans, dan lapisan anotasi. Pada kondisi utama, hanya citra RGB depan, kedalaman terdaftar, keadaan UAV, ringkasan kemajuan misi, dan instruksi yang memasuki \texttt{CanonicalObservation}. Segmentasi instans, label objek, dan transformasi objek merupakan data oracle yang hanya tersedia bagi pembangun manifest, \textit{baseline} oracle, dan evaluator; adaptor, pengelola misi, gerbang keselamatan, serta pengendali tidak dapat membacanya untuk memilih aksi. % [SRC-B2-12]

ROS~2 menyediakan middleware robotika yang berpusat pada simpul, topik, layanan, aksi, parameter, dan transformasi \citep{macenski2022ros2}. Jembatan penelitian menerbitkan citra kamera, \textit{CameraInfo}, odometri, IMU, GPS, data jarak, dan transformasi objek; jembatan juga menerima perintah kecepatan serta tujuan posisi lokal. Perbedaan kerangka koordinat antara simulator, ROS~2, dan badan UAV harus dikelola pada pembangun pengamatan serta adaptor pengendali. Adaptor VLM tidak boleh melakukan konversi koordinat secara mandiri karena konversi ganda dapat menghasilkan tindakan yang tidak konsisten dengan keadaan kendaraan. % [SRC-B2-06]

\subsection{Adaptor Agnostik-Model dan Aksi Terstruktur}
\label{subsec:adapter_aksi}
Adaptor agnostik-model menyembunyikan perbedaan permintaan dan respons setiap \textit{backend} di balik antarmuka yang tetap. Pada DroneVL, adaptor mendeklarasikan kapabilitas, menyiapkan permintaan dari \texttt{CanonicalObservation}, menjalankan inferensi, lalu mengembalikan respons mentah beserta metadata. Pengurai memiliki tanggung jawab tunggal untuk mendeserialisasi respons khusus-\textit{backend} menjadi objek aksi kanonis; kegagalan pada tahap ini dicatat sebagai kegagalan penguraian, bukan kegagalan adaptor. Adaptor tidak berlangganan topik ROS secara langsung, tidak mengendalikan UAV, tidak menentukan keberhasilan misi, dan tidak dapat melewati gerbang keselamatan. Batas tanggung jawab ini membuat penggantian \textit{backend} dapat dilakukan tanpa mengubah prosedur pengambilan sensor, pengendalian, maupun evaluasi.

Kontrak DroneVL membedakan perintah kendali dari usulan peristiwa misi. Perintah kendali yang dapat diteruskan menuju pengendali tingkat rendah adalah \texttt{MOVE\_RELATIVE}, untuk gerakan relatif berbatas; \texttt{ROTATE\_YAW}, untuk perubahan orientasi yaw; \texttt{HOVER\_AND\_OBSERVE}, untuk memperoleh pengamatan baru; dan \texttt{ABORT}, untuk penghentian terkendali. \texttt{DECLARE\_TARGET} serta \texttt{ADVANCE\_SUBGOAL} adalah usulan semantik kepada pengelola misi, bukan perintah penerbangan. Pengelola misi memverifikasi usulan tersebut terhadap predikat episode dan \textit{ground truth} evaluator sebelum mengubah keadaan misi.

Pengurai memeriksa struktur respons, bidang wajib, tipe data, nilai numerik berhingga, rentang tindakan, dan legalitas aksi pada keadaan misi saat ini. Satu \textit{prompt} perbaikan dapat digunakan apabila respons tidak valid; setelah $N_{\mathrm{parse}}$ kegagalan penguraian berturut-turut yang nilainya dibekukan pada manifest episode, pengelola misi menghentikan episode dengan alasan terminal \texttt{PARSE\_FAILURE}. Penghitung direset hanya setelah aksi kanonis yang valid diterima. Teks penjelasan tetap dicatat sebagai bukti, tetapi tidak digunakan sebagai parameter penerbangan.

\subsection{Batasan Keselamatan dan \textit{Benchmark} yang Dapat Direproduksi}
\label{subsec:keselamatan_benchmark}
Gerbang keselamatan pada DroneVL adalah mekanisme operasional untuk simulasi, bukan sertifikasi keselamatan penerbangan nyata. Gerbang tersebut memeriksa validitas skema, batas gerak dan yaw, geofence, ketinggian, usia data, jarak bebas dari informasi kedalaman atau rintangan, serta anggaran waktu dan keputusan. Setiap aksi dapat diterima, dipotong, diganti, atau ditolak dengan alasan yang dicatat. Log menyimpan aksi kanonis yang diusulkan, keputusan gerbang, aksi aktual, dan alasan intervensi agar keselamatan kebijakan mentah dapat dibedakan dari keselamatan sistem yang dibantu.

\textit{Benchmark} \textit{closed-loop} yang dapat direproduksi harus mendefinisikan lebih dari kumpulan gambar. Setiap episode memuat adegan, \textit{seed}, pose awal, instruksi, target, kondisi lingkungan, anggaran waktu dan keputusan, batasan, kondisi terminal, dan \textit{ground truth}. Manifest berversi juga memuat label kemampuan yang dapat bertumpang tindih, pemisahan pengembangan dan pengujian, serta versi kode, aset, model, \textit{prompt}, perangkat keras, dan log mentah. Split pengujian dibekukan sebelum penyetelan LoRA/GRPO dan dipisahkan dari pengembangan berdasarkan tata ruang, instans objek, komposisi instruksi, serta kondisi visual. Kebutuhan evaluasi kecerdasan spasial yang mempertimbangkan karakteristik navigasi UAV ditekankan oleh \textit{benchmark} VLM untuk UAV \citep{zhang2025skyready}. DroneVL Arena menghitung metrik dari log evaluator, bukan dari tabel yang diedit secara manual. % [SRC-B2-07]

%-----------------------------------------------------------------------------%
\section{Penelitian Terkait}
\label{sec:penelitian_terkait}
%-----------------------------------------------------------------------------%

\subsection{\textit{Vision-Language Model}, \textit{Embodied Agent}, dan Navigasi UAV}
\textit{Vision-language model} menunjukkan bahwa representasi visual dan bahasa dapat digunakan untuk memahami objek dan instruksi, tetapi representasi tersebut belum menentukan sendiri batas tindakan yang aman bagi robot \citep{radford2021clip,vinyals2015show}. Pada konteks UAV, van der Wal mengeksplorasi penggunaan VLM untuk eksplorasi aktif pada inspeksi anomali, yang menunjukkan relevansi penalaran visual--bahasa bagi tugas pengamatan dari udara \citep{wal2026active}. DroneVL mengambil posisi berbeda dengan menempatkan VLM di balik kontrak aksi yang terbatas, pengurai, dan gerbang keselamatan. Dengan demikian, evaluasi tidak hanya menilai apakah model menghasilkan respons semantik yang masuk akal, tetapi juga apakah respons tersebut dapat dinormalisasi dan dieksekusi dalam lingkar kendali yang tetap. % [SRC-B2-08]

Penelitian \textit{vision-language navigation} menekankan pelaksanaan instruksi dari pengamatan berurutan, penambatan visual, dan pemantauan kemajuan \citep{zhang2022explicit,ma2019selfmonitoring}. Tinjauan navigasi tujuan-objek pada UAV juga menunjukkan ketergantungan antara persepsi, pelokalan, perencanaan, dan kendali serta kebutuhan standardisasi kerangka evaluasi \citep{ayala2024uav}. \textit{Benchmark} \textit{Is your VLM Sky-Ready?} secara khusus mengarahkan evaluasi pada kecerdasan spasial untuk navigasi UAV \citep{zhang2025skyready}. DroneVL Arena sejalan dengan arah tersebut melalui tugas penambatan, penalaran spasial, pencarian, navigasi, dan instruksi multilangkah; pembeda utamanya adalah penggunaan satu kontrak pengamatan, aksi, pengendali, dan prosedur pencatatan untuk membandingkan \textit{backend} yang berbeda. % [SRC-B2-09]

\subsection{VLA Khusus UAV: CognitiveDrone}
\label{subsec:cognitivedrone_terkait}
CognitiveDrone merupakan penelitian yang secara langsung menerapkan \textit{vision-language-action} pada UAV. Lykov dkk. memperkenalkan model yang dilatih pada lebih dari 8.000 trajektori penerbangan simulasi untuk tugas pengenalan manusia, pemahaman simbol, dan penalaran; model menerima masukan visual orang-pertama serta instruksi teks, kemudian menghasilkan perintah aksi empat-dimensi secara waktu nyata \citep{lykov2025cognitivedrone}. Penelitian tersebut juga memperkenalkan CognitiveDroneBench sebagai \textit{benchmark} untuk tugas kognitif UAV dan varian CognitiveDrone-R1 yang menambahkan modul penalaran VLM sebelum kendali berfrekuensi tinggi.

Relevansi CognitiveDrone bagi penelitian ini terletak pada dua aspek. Pertama, ia menyediakan titik pembanding VLA yang asumsi perwujudannya sudah dekat dengan UAV. Kedua, ia menunjukkan bahwa evaluasi tindakan kognitif UAV memerlukan lebih dari pengenalan visual statis. Namun, hasil CognitiveDroneBench tidak diperlakukan sebagai hasil yang dapat dibandingkan langsung dengan DroneVL Arena karena lingkungan, episode, skema aksi, dan prosedur evaluasinya berbeda. Dalam DroneVL, CognitiveDrone ditempatkan sebagai satu kondisi pada eksperimen portabilitas VLA; keluarannya tetap harus menyatakan skema aksi, satuan, kerangka koordinat, serta horizon sebelum dapat dinormalisasi menjadi aksi kanonis.

Dengan posisi tersebut, DroneVL melengkapi, bukan menggantikan, CognitiveDrone. CognitiveDrone menilai kebijakan VLA khusus UAV dalam \textit{benchmark}-nya sendiri, sedangkan DroneVL menguji apakah kebijakan khusus UAV tersebut, OpenVLA-7B yang diadaptasi untuk UAV, dan OpenPI $\pi_{0.5}$ yang diadaptasi untuk UAV dapat melewati kontrak pengamatan, penguraian, gerbang keselamatan, dan evaluator yang sama. Pemisahan ini memungkinkan kegagalan karena format keluaran atau asumsi perwujudan dibedakan dari kegagalan navigasi pada setiap model.

\subsection{Celah Integrasi, Pengaman, dan \textit{Benchmark}}
CoSyS-AirSim dan ROS~2 menyediakan dua fondasi yang saling melengkapi, tetapi berada pada lapisan yang berbeda. CoSyS-AirSim memperluas AirSim melalui sensor, jenis kendaraan, dan lingkungan yang dapat diubah untuk simulasi industri serta pengujian algoritme navigasi \citep{jansen2023cosysairsim}. ROS~2 menyediakan arsitektur middleware berbasis simpul, topik, layanan, aksi, parameter, dan transformasi untuk menghubungkan komponen robotika \citep{macenski2022ros2}. Kedua karya tersebut menjadi dasar yang kuat untuk sensor, simulasi, dan komunikasi pada penelitian ini, tetapi kontribusi yang dilaporkan berfokus pada infrastruktur, bukan pada kontrak adaptor yang menormalkan keluaran VLM dan VLA heterogen menjadi aksi navigasi UAV yang dapat diaudit.

Pada lapisan model dan evaluasi, penelitian terkait juga menutup bagian masalah yang berbeda. CognitiveDrone menyediakan VLA khusus UAV beserta CognitiveDroneBench, dengan keluaran aksi empat-dimensi dari pengamatan orang-pertama dan instruksi \citep{lykov2025cognitivedrone}. Sementara itu, SpatialSky-Bench dalam \textit{Is your VLM Sky-Ready?} mengevaluasi kecerdasan spasial VLM untuk navigasi UAV melalui kategori persepsi lingkungan dan pemahaman adegan yang dibagi menjadi 13 subkategori \citep{zhang2025skyready}. Kedua karya tersebut menunjukkan bahwa VLA khusus UAV dan evaluasi kemampuan spasial VLM sama-sama diperlukan, tetapi memakai model, tugas, representasi keluaran, dan prosedur evaluasi yang dirancang untuk tujuan masing-masing.

Berdasarkan sintesis tersebut, celah yang ditangani DroneVL adalah lapisan integrasi dan evaluasi yang memungkinkan beberapa VLM serta VLA diuji pada kontrak pengamatan, aksi kanonis, pengendali tingkat rendah, batas keselamatan, episode, dan pencatatan yang sama. Pernyataan ini merupakan posisi penelitian yang diturunkan dari cakupan karya-karya di atas, bukan klaim bahwa karya terdahulu tidak memiliki nilai integrasi atau evaluasi. Celah tersebut diuji dengan melaporkan setiap tahap, yaitu inferensi, penguraian, keputusan gerbang keselamatan, aksi aktual, dan hasil misi, sehingga perbedaan kegagalan antarmodel dapat ditelusuri, bukan hanya dibandingkan melalui satu angka keberhasilan.

%-----------------------------------------------------------------------------%
\section{Posisi Penelitian}
\label{sec:posisi_penelitian}
%-----------------------------------------------------------------------------%
Tabel~\ref{tab:posisi_penelitian} merangkum posisi DroneVL terhadap kelompok penelitian terkait. Kontribusi penelitian bukan VLM baru atau klaim kelayakan penerbangan nyata, melainkan arsitektur integrasi dan evaluasi yang membuat perbandingan \textit{backend} dapat ditelusuri pada lingkungan UAV simulasi.

\begin{table}[h!]
    \centering
    \captionsetup{justification=justified, singlelinecheck=false}
    \caption{Posisi DroneVL terhadap kelompok penelitian terkait.}
    \label{tab:posisi_penelitian}
    \small
    \begin{tabularx}{\linewidth}{|>{\raggedright\arraybackslash}p{3.0cm}|>{\raggedright\arraybackslash}X|>{\raggedright\arraybackslash}X|>{\raggedright\arraybackslash}X|}
        \hline
        \textbf{Aspek} & \textbf{VLM dan navigasi berwujud} & \textbf{CoSyS-AirSim / ROS~2} & \textbf{DroneVL} \\
        \hline
        Pemahaman visual--bahasa & Tersedia dan bergantung pada model. & Hanya fondasi sensor. & Tersedia melalui \textit{backend} VLM. \\
        \hline
        Domain utama & Gambar statis atau agen darat pada banyak penelitian. & Simulasi kendaraan dan robot. & UAV multirotor simulasi. \\
        \hline
        Navigasi UAV \textit{closed-loop} & Tidak selalu menjadi fokus. & Menyediakan fondasi kendali. & Menjadi konfigurasi evaluasi utama. \\
        \hline
        \textit{Backend} VLM yang dapat diganti & Umumnya tidak dibahas. & Tidak disediakan. & Disediakan melalui adaptor. \\
        \hline
        Kontrak aksi dan gerbang keselamatan & Bergantung pada implementasi. & Kendali dasar tersedia. & Ditetapkan bersama dan dicatat. \\
        \hline
        Episode semantik dan evaluator berversi & Bergantung pada \textit{benchmark}. & Tidak disediakan sebagai arena semantik. & Disediakan oleh DroneVL Arena. \\
        \hline
    \end{tabularx}
\end{table}

Secara ringkas, DroneVL memiliki tiga kontribusi yang saling terkait. Pertama, penelitian ini merancang arsitektur agnostik-model yang memisahkan simulator, adaptor, pengurai, gerbang keselamatan, dan pengendali tingkat rendah. Kedua, DroneVL Arena mendefinisikan episode berversi, anotasi, pemisahan data, log, dan evaluator untuk tugas navigasi semantik UAV. Ketiga, penelitian ini menyiapkan perbandingan empiris yang melaporkan kualitas, keselamatan, efisiensi, aksi tidak valid, penggunaan sumber daya, dan latensi. Batas kontribusi tersebut adalah lingkungan simulasi dan kondisi eksperimen yang ditetapkan; penelitian ini tidak mensertifikasi keselamatan penerbangan nyata.

%-----------------------------------------------------------------------------%
\chapter{\babTiga}
\label{cha:metodologi}
%-----------------------------------------------------------------------------%
Bab ini menjelaskan metodologi penelitian dan rancangan sistem DroneVL. Rancangan memisahkan inferensi VLM atau VLA yang bergantung pada \textit{backend} dari komponen navigasi yang dibekukan pada seluruh kondisi eksperimen. Spesifikasi berikut menjadi acuan implementasi, pencatatan, dan evaluasi; setiap nilai konfigurasi yang belum bersifat universal dibekukan dalam manifest eksperimen sebelum pengujian dimulai.

%-----------------------------------------------------------------------------%
\section{Desain Penelitian}
\label{sec:desain_penelitian}
%-----------------------------------------------------------------------------%
Penelitian menggunakan pendekatan rekayasa sistem dan evaluasi eksperimental. Artefak yang dibangun meliputi DroneVL sebagai adaptor untuk VLM dan VLA, serta DroneVL Arena sebagai kerangka evaluasi \textit{closed-loop}. Evaluasi membedakan dua kondisi. Kondisi antarmuka bersama menyetarakan episode, sensor yang diizinkan, resolusi masukan, riwayat, \textit{prompt}, frekuensi keputusan, batas waktu, kebijakan coba-ulang, pengendali, dan evaluator sejauh didukung semua \textit{backend}. Kondisi kapabilitas asli mencatat konfigurasi terbaik yang sah bagi setiap \textit{backend}; hasilnya dilaporkan terpisah dan tidak ditafsirkan sebagai pengaruh model semata.

Eksperimen dibagi menjadi tiga keluarga. Eksperimen VLM membandingkan GPT-5.6, Qwen3.6 35B-A3B, dan Gemini sebagai perencana semantik. Eksperimen portabilitas VLA membandingkan CognitiveDrone, OpenVLA-7B yang diadaptasi untuk UAV, dan OpenPI $\pi_{0.5}$ yang diadaptasi untuk UAV. Eksperimen utama adaptasi membandingkan model dasar OpenVLA-7B (M0), OpenVLA-7B dengan LoRA yang diadaptasi untuk UAV (M1), dan M1 yang dilanjutkan GRPO (M2). Hasil ketiga keluarga dilaporkan terpisah; perbandingan lintas keluarga hanya bersifat deskriptif karena representasi aksi dan riwayat pelatihan berbeda.

Setiap episode memiliki pasangan percobaan yang sama di dalam keluarga eksperimen yang relevan. Kesamaan jalur hilir mengendalikan perubahan pengendali dan skenario, tetapi tidak menghapus perbedaan bawaan seperti modalitas, panjang konteks, atau latensi layanan. Karena itu, konfigurasi efektif dan kehilangan informasi akibat normalisasi dicatat bersama hasil eksperimen.

%-----------------------------------------------------------------------------%
\section{Rancangan Arsitektur DroneVL}
\label{sec:rancangan_dronevl}
%-----------------------------------------------------------------------------%
Arsitektur terdiri atas komponen pembentuk pengamatan, adaptor \textit{backend}, pengurai respons, pengelola misi, gerbang keselamatan, pengendali tingkat rendah, dan pencatat eksperimen. Alur representasi dibekukan sebagai berikut: \textit{CanonicalObservation} $\rightarrow$ permintaan khusus-\textit{backend} $\rightarrow$ respons mentah dan metadata $\rightarrow$ hasil deserialisasi khusus-\textit{backend} $\rightarrow$ aksi kanonis atau usulan peristiwa misi $\rightarrow$ keputusan keselamatan $\rightarrow$ aksi aktual. Adaptor hanya bertanggung jawab atas dua tahap khusus-\textit{backend}; pengurai memiliki satu tanggung jawab untuk membentuk representasi kanonis dan melaporkan kegagalan penguraian.

Pada jalur eksperimen VLM, GPT-5.6, Qwen3.6 35B-A3B, dan Gemini hanya terhubung ke jalur penerbangan melalui DroneVL Model Adapter. GPT-5.6 menggunakan antarmuka API gambar--teks, Qwen3.6 35B-A3B menggunakan antarmuka gambar--teks dengan arsitektur \textit{mixture-of-experts} (MoE) multimodal, sedangkan Gemini menggunakan Live API dua arah untuk menerima konteks multimodal dan mengembalikan teks atau pemanggilan fungsi. Respons \textit{streaming} tersebut tidak diteruskan langsung ke UAV; respons harus dihimpun pada batas keputusan, dideserialisasi, lalu dinormalisasi menjadi aksi tingkat tinggi yang dibatasi atau ditolak sebelum memasuki gerbang keselamatan.

Jalur VLA menggunakan kontrak pengamatan kanonis yang sama, tetapi adaptor VLA juga membaca deskriptor perwujudan. Deskriptor tersebut memuat skema aksi, satuan, kerangka koordinat, horizon, dan versi \textit{checkpoint}. CognitiveDrone, OpenVLA-7B yang diadaptasi untuk UAV, dan OpenPI $\pi_{0.5}$ yang diadaptasi untuk UAV mengembalikan \textit{NativeActionEnvelope} berversi. Adaptor mendeserialisasi amplop tersebut ke kandidat aksi kanonis, tanpa meneruskan aksi asli langsung ke UAV.

OpenVLA-7B menggunakan token aksi UAV yang dipelajari dari split pelatihan. OpenPI $\pi_{0.5}$ menggunakan potongan aksi UAV kontinu yang didekode per horizon, sedangkan CognitiveDrone menggunakan skema aksi empat-dimensi yang ditetapkan dalam manifest. Semua kandidat selanjutnya melewati pengurai dan gerbang keselamatan yang sama.

Perbedaan bentuk masukan dan keluaran ketiga model sebelum normalisasi menjadi objek pencatatan pada lapisan adaptor, bukan alasan untuk mengubah pengendali, evaluator, atau aturan terminal.

\subsection{Kontrak Pengamatan \textit{CanonicalObservation}}
\label{subsec:canonical_observation}
Setiap keputusan menerima \texttt{CanonicalObservation} berversi yang sedikitnya memuat \texttt{schema\_version}, \texttt{episode\_id}, \texttt{decision\_id}, waktu simulator, waktu dinding, dan cap waktu penangkapan sensor. Isinya terdiri atas citra RGB depan; kedalaman terdaftar beserta satuan dan kerangka koordinat kamera; keadaan UAV yang memuat pose, kecepatan, dan orientasi; instruksi; serta ringkasan kemajuan misi yang dihasilkan pengelola misi. Manifest menentukan resolusi, pemotongan, normalisasi, usia maksimum frame, frekuensi keputusan, dan cara riwayat dipadatkan.

Segmentasi instans, label objek, transformasi objek, serta anotasi target tidak berada dalam \texttt{CanonicalObservation}. Data tersebut hanya dapat dibaca pembangun manifest, \textit{baseline} oracle, dan evaluator melalui kanal oracle terpisah. Pembangun pengamatan mengunci pasangan RGB, kedalaman, dan odometri pada waktu simulator yang sama; frame yang melampaui usia maksimum dibuang, bukan dikirimkan ke \textit{backend}.

\subsection{Kontrak Aksi dan Peristiwa Misi}
\label{subsec:canonical_action}
Aksi kendali kanonis memuat \texttt{action\_type}, \texttt{observation\_id}, \texttt{sequence\_id}, \texttt{valid\_until}, dan parameter bernilai hingga. \texttt{MOVE\_RELATIVE} memuat perpindahan $(\Delta x,\Delta y,\Delta z)$ dalam meter pada kerangka badan UAV yang dinyatakan manifest; \texttt{ROTATE\_YAW} memuat perubahan yaw dalam radian; \texttt{HOVER\_AND\_OBSERVE} memuat durasi terbatas; dan \texttt{ABORT} tidak memiliki parameter gerak. Satu konversi kerangka dari aksi kanonis ke perintah pengendali dilakukan pada batas pengendali, sehingga adaptor tidak dapat melakukan transformasi koordinat tambahan.

\texttt{DECLARE\_TARGET} dan \texttt{ADVANCE\_SUBGOAL} bukan aksi kendali. Keduanya adalah usulan peristiwa misi yang memuat identitas atau predikat yang diklaim; pengelola misi memverifikasi usulan terhadap kondisi episode dan kanal oracle evaluator sebelum mengubah keadaan misi. Penolakan usulan dicatat sebagai kegagalan semantik, bukan sebagai kegagalan pengendali.

Pengurai memeriksa skema, bidang wajib, tipe data, nilai numerik berhingga, rentang aksi, cap waktu, dan legalitas aksi pada keadaan misi saat ini. Satu \textit{prompt} perbaikan dapat digunakan untuk respons tidak valid. Setelah $N_{\mathrm{parse}}$ kegagalan penguraian berturut-turut, dengan nilai yang dibekukan dalam manifest, episode diakhiri dengan \texttt{PARSE\_FAILURE}; penghitung hanya direset oleh aksi kanonis yang valid.

Untuk VLA, \textit{NativeActionEnvelope} wajib memuat \texttt{action\_schema\_version}, \texttt{checkpoint\_id}, \texttt{native\_action}, \texttt{coordinate\_frame}, \texttt{unit}, dan \texttt{horizon}. Pemetaan ke \texttt{MOVE\_RELATIVE}, \texttt{ROTATE\_YAW}, \texttt{HOVER\_AND\_OBSERVE}, atau \texttt{ABORT} ditetapkan per \textit{backend} dalam manifest dan diuji pada split pengembangan. Tidak ada pemetaan implisit dari aksi sendi robot, gripper, atau ruang aksi manipulasi ke aksi UAV. Apabila skema, \textit{checkpoint}, atau pemetaan tidak cocok, respons ditolak sebagai \texttt{VLA\_SCHEMA\_FAILURE} dan tidak dieksekusi.

%-----------------------------------------------------------------------------%
\section{Rancangan Lingkungan DroneVL Arena}
\label{sec:rancangan_arena}
%-----------------------------------------------------------------------------%
Setiap episode DroneVL Arena didefinisikan oleh manifest berversi yang memuat identitas dan versi adegan, \textit{seed}, pose awal, instruksi, target, kondisi lingkungan, anggaran waktu dan keputusan, batasan keselamatan, kondisi terminal, serta anotasi \textit{ground truth}. Setiap episode menerima satu atau lebih label kemampuan: penambatan visual, penalaran spasial, pencarian objek, navigasi semantik, dan instruksi multilangkah. Hasil dilaporkan menurut hasil misi akhir dan menurut label kemampuan; satu episode multikemampuan tidak dipaksa masuk ke satu kategori eksklusif.

Manifest membedakan split pelatihan, pengembangan, dan pengujian. Split pengujian dibekukan sebelum penyetelan dan dipisahkan dari pelatihan maupun pengembangan berdasarkan tata ruang, instans objek, komposisi instruksi, dan kondisi visual. Keberhasilan misi memerlukan pemenuhan target, predikat relasi bila ada, urutan subtujuan, toleransi pose, dan durasi bertahan yang ditetapkan oleh manifest. Untuk relasi spasial, evaluator memeriksa kecocokan instans target dan predikat pose secara terpisah dengan kerangka acuan, toleransi, serta waktu penilaian yang eksplisit.

%-----------------------------------------------------------------------------%
\section{Protokol Eksperimen, Log, dan Metrik}
\label{sec:protokol_eksperimen}
%-----------------------------------------------------------------------------%
Sebelum pengujian VLM, konfigurasi GPT-5.6, Qwen3.6 35B-A3B dengan tag \texttt{qwen3.6:35b-a3b}, dan Gemini dikunci serta dicatat. Untuk Gemini, catatan juga memuat mode sesi Live API, modalitas \textit{streaming}, batas waktu pengumpulan respons pada setiap keputusan, serta aturan pemetaan pemanggilan fungsi atau teks ke respons mentah.

Sebelum pengujian portabilitas VLA, konfigurasi CognitiveDrone, OpenVLA-7B yang diadaptasi untuk UAV, dan OpenPI $\pi_{0.5}$ yang diadaptasi untuk UAV dikunci serta dicatat. Catatan tersebut memuat asal \textit{checkpoint}, split dan versi trajektori adaptasi, skema atau tokenisasi aksi UAV, horizon keluaran, serta prosedur dekode. Untuk seluruh \textit{backend}, catatan juga memuat identitas versi atau rilis, penyedia dan jenis antarmuka, kapabilitas yang dideklarasikan, prapemrosesan gambar, resolusi efektif, \textit{prompt}, parameter dekode, \textit{runtime}, perangkat keras lokal, dan wilayah atau konfigurasi layanan apabila relevan. Seluruh \textit{backend} pada kondisi antarmuka bersama menggunakan episode, pengendali, batasan keselamatan, kebijakan coba-ulang, dan evaluator yang sama.

Log per keputusan memuat identitas episode dan manifest, versi anotasi, \texttt{observation\_id}, respons mentah, hasil deserialisasi, aksi kanonis yang diusulkan, usulan peristiwa misi, keputusan gerbang, aksi aktual, keadaan UAV, usia pengamatan, waktu inferensi, waktu keputusan ujung-ke-ujung, dan alasan terminal. Untuk penambatan, log juga menyimpan target prediksi, target acuan, aturan pencocokan, hasil kecocokan instans, dan hasil predikat relasi. Dengan demikian, metrik agregat dapat dihitung ulang dari log mentah.

Akurasi penambatan instans didefinisikan sebagai $A_{\mathrm{inst}}=N^{-1}\sum_{i=1}^{N}\mathbf{1}[p_i=g_i]$ pada semua peluang deklarasi target yang dievaluasi. Akurasi relasi dihitung terpisah pada peluang yang memuat predikat spasial. Tingkat keberhasilan misi adalah proporsi episode yang memenuhi seluruh kondisi terminal berhasil. Tingkat tabrakan adalah proporsi episode yang mengalami sedikitnya satu tabrakan; jumlah tabrakan per keputusan dilaporkan sebagai metrik tambahan. Tingkat aksi tidak valid adalah proporsi respons awal yang gagal diurai, sedangkan kegagalan setelah \textit{prompt} perbaikan dicatat terpisah. Tingkat intervensi keselamatan adalah proporsi aksi kanonis yang dipotong, diganti, atau ditolak. Efisiensi lintasan adalah rasio panjang lintasan referensi yang layak terhadap panjang lintasan aktual untuk episode berhasil, dengan lintasan referensi dibekukan oleh manifest. Latensi inferensi dan latensi ujung-ke-ujung dilaporkan sebagai median, persentil ke-95, dan distribusi per episode.

%-----------------------------------------------------------------------------%
\section{Eksperimen Adaptasi LoRA dan GRPO}
\label{sec:lora_grpo_experiment}
%-----------------------------------------------------------------------------%
LoRA dan GRPO merupakan tujuan utama penelitian pada OpenVLA-7B yang diadaptasi untuk UAV. Tiga kondisi dibandingkan: model dasar OpenVLA-7B dengan tokenisasi aksi UAV yang ditetapkan (M0); M0 dengan LoRA terawasi menggunakan trajektori CoSyS-AirSim pada split pelatihan (M1); dan M1 yang dilanjutkan GRPO menggunakan \textit{rollout} CoSyS-AirSim (M2). Data pelatihan hanya berasal dari split pelatihan; split pengembangan digunakan untuk memilih hiperparameter; split pengujian yang telah dibekukan tidak digunakan untuk penyetelan maupun \textit{rollout}.

LoRA dan GRPO dibatasi pada VLA OpenVLA-7B, bukan diterapkan pada seluruh model dalam penelitian. Tujuannya adalah mengisolasi pengaruh adaptasi kebijakan aksi UAV melalui perbandingan M0--M1--M2 dengan data, pengendali, episode, dan metrik yang sama. OpenVLA-7B dipilih karena keluarannya berupa token aksi diskret dan prosedur LoRA tersedia untuk tugas maupun perwujudan robot baru \citep{kim2024openvla}; kondisi ini memungkinkan log-probabilitas token aksi dicatat pada setiap \textit{rollout} untuk perhitungan GRPO. Kelompok VLM tetap dievaluasi sebagai perencana semantik melalui kontrak pengamatan dan aksi yang sama. Penerapan LoRA atau GRPO pada VLM akan menambahkan pertanyaan penelitian lain tentang adaptasi perencana, sekaligus mengaburkan apakah perubahan hasil berasal dari model perencana, kebijakan aksi, atau lapisan adaptor.

CognitiveDrone dan OpenPI $\pi_{0.5}$ tetap digunakan pada eksperimen portabilitas VLA, tetapi bukan sebagai kondisi LoRA+GRPO. CognitiveDrone menjadi pembanding kebijakan khusus UAV yang sudah tersedia. OpenPI $\pi_{0.5}$ dapat diadaptasi terhadap trajektori UAV sesuai prosedur pelatihannya sendiri, tetapi keluaran \textit{flow-matching}-nya berupa potongan aksi kontinu \citep{physicalintelligence2025openpi}. Karena representasi tersebut tidak serta-merta menyediakan rasio probabilitas kebijakan yang diperlukan GRPO, memasukkannya ke M0--M1--M2 akan menghasilkan perbandingan metode optimisasi yang tidak sepadan. Dengan pembatasan ini, eksperimen portabilitas menjawab kesesuaian antarmuka VLA, sedangkan eksperimen OpenVLA-7B menjawab dampak LoRA dan GRPO secara terisolasi.

Manifest pelatihan menetapkan data dan versi \textit{prompt}, tokenisasi serta skema aksi UAV, lapisan dan peringkat LoRA, kuantisasi, perangkat keras, anggaran token, jumlah \textit{rollout}, ukuran kelompok GRPO, kebijakan acuan, dan \textit{seed}. Pada M2, sejumlah \textit{rollout} bertitik awal sama disampel dari token aksi OpenVLA, log-probabilitas kebijakan lama disimpan, dan penghargaan relatif kelompok dihitung dari keberhasilan misi, kemajuan subtujuan, keselamatan, kepatuhan format aksi, serta penalti aksi tidak valid. Bobot komponen serta setiap ablasi dibekukan sebelum pelatihan. Ketiga kondisi dievaluasi dengan manifest episode, pengendali, batas keselamatan, dan metrik yang sama, sehingga perubahan hasil dapat ditafsirkan sebagai efek LoRA dan GRPO pada OpenVLA-7B UAV yang diuji. OpenPI $\pi_{0.5}$ tidak dimasukkan ke klaim GRPO; ia hanya dibandingkan pada eksperimen portabilitas VLA setelah adaptasi UAVnya sendiri.

%-----------------------------------------------------------------------------%
\chapter{\babEmpat}
\label{cha:pengujian}
%-----------------------------------------------------------------------------%
Bab ini akan menyajikan realisasi implementasi, hasil pengujian DroneVLM-Arena, dan analisis perbandingan \textit{backend} VLM. Seluruh hasil akan dilaporkan berdasarkan log eksperimen yang dihasilkan oleh evaluator.

%-----------------------------------------------------------------------------%
\section{Realisasi Eksperimen dan Pengolahan Data}
\label{sec:realisasi_eksperimen}
%-----------------------------------------------------------------------------%
Bagian ini akan memaparkan konfigurasi simulator, implementasi adaptor, konfigurasi model, episode yang digunakan, serta prosedur pengolahan log eksperimen.

%-----------------------------------------------------------------------------%
\section{Hasil dan Analisis}
\label{sec:hasil_analisis}
%-----------------------------------------------------------------------------%
Bagian ini akan menyajikan hasil penambatan visual, penalaran spasial, pencarian objek, navigasi semantik, dan instruksi multilangkah. Analisis akan membahas keberhasilan misi, keselamatan, efisiensi lintasan, keluaran tidak valid, penggunaan sumber daya, dan latensi.

%-----------------------------------------------------------------------------%
\chapter{\babLima}
\label{cha:penutup}
%-----------------------------------------------------------------------------%

%-----------------------------------------------------------------------------%
\section{Kesimpulan}
\label{sec:kesimpulan}
%-----------------------------------------------------------------------------%
Bagian ini akan memuat kesimpulan yang menjawab rumusan masalah berdasarkan hasil eksperimen dan analisis pada Bab~IV.

%-----------------------------------------------------------------------------%
\section{Saran}
\label{sec:saran}
%-----------------------------------------------------------------------------%
Bagian ini akan memuat saran untuk pengembangan DroneVL, perluasan DroneVL Arena, evaluasi model yang lebih beragam, dan penelitian \textit{sim-to-real} pada masa mendatang.


%
% Daftar Pustaka
%\input{pustaka}

% Alternatif manajemen daftar pustaka dengan \bibliography
\bibliographystyle{IEEEtran}
\bibliography{pustaka}

%
% Lampiran 
%
%\begin{appendix}
%    \input{_internals/markLampiran}
%    \setcounter{page}{2}\textbf{}
%    %---------------------------------------------------------------------%
\addChapter{Lampiran 1 - Kode Program}
\chapter*{Lampiran 1 - Kode Program}
%---------------------------------------------------------------------%

\section*{Lampiran 1.1 Contoh Penulisan Kode Program (\texttt{contoh.py})}
\addcontentsline{toc}{section}{Lampiran 1.1 Contoh Penulisan Kode Program (contoh.py)}

Lampiran digunakan untuk memuat berkas pendukung seperti kode program, data
mentah, atau dokumen tambahan. Potongan kode ditulis menggunakan lingkungan
\texttt{lstlisting} dengan gaya \texttt{mystyle} yang telah didefinisikan pada
\texttt{thesis.tex}.

\begin{lstlisting}[
    style=mystyle,
    language=Python,
]
def main():
    """Ganti isi lampiran ini dengan kode program yang sebenarnya."""
    data = [1, 2, 3, 4, 5]
    total = sum(data)
    print("Jumlah data:", len(data))
    print("Total nilai:", total)


if __name__ == "__main__":
    main()
\end{lstlisting}

%---------------------------------------------------------------------%
\section*{Lampiran 1.2 Contoh Lampiran Tabel}
\addcontentsline{toc}{section}{Lampiran 1.2 Contoh Lampiran Tabel}

Tabel pendukung yang terlalu panjang untuk dimuat pada batang tubuh laporan
dapat diletakkan pada lampiran.

\begin{table}[htbp]
    \centering
    \addtocounter{table}{-1}
    \caption{Contoh tabel pada lampiran}
    \label{tab:contoh_lampiran}
    \small
    \setlength{\tabcolsep}{4pt}
    \renewcommand{\arraystretch}{1.15}
    \begin{tabularx}{\textwidth}{
        >{\centering\arraybackslash}p{2cm}
        >{\raggedright\arraybackslash}p{3cm}
        >{\raggedright\arraybackslash}X
    }
        \toprule
        \textbf{No.} & \textbf{Parameter} & \textbf{Keterangan} \\
        \midrule
        1 & Parameter A & Keterangan parameter pertama. \\
        2 & Parameter B & Keterangan parameter kedua. \\
        3 & Parameter C & Keterangan parameter ketiga. \\
        \bottomrule
    \end{tabularx}
\end{table}

%\end{appendix}

\end{document}
